\begin{titlepage}
\centering
\vspace*{2.2cm}

{\LARGE\bfseries Mytheus\par}
\vspace{0.5em}
{\large Personalized Multi-Agent Meeting Orchestration\par}

\vspace{1.6cm}
{\color{rulegrey}\hrule}
\vspace{1.1cm}

{\Large\bfseries Research Gaps and Novel Contributions\par}
\vspace{0.5em}
{\large What twenty five papers leave open,\\ and what this project can defensibly claim\par}

\vspace{1.1cm}
{\color{rulegrey}\hrule}
\vspace{2.0cm}

\begin{minipage}{0.82\textwidth}
\small
This is the companion to \textbf{A Literature Foundation}. That document
established what twenty five verified papers contribute. This one asks the
harder question: what none of them do.

\vspace{0.7em}
The limitations recorded at the foot of every entry in the companion, one
hundred and thirty two of them, were read back through six analytical lenses.
The recurring gaps became candidate contributions, and every candidate was
then handed to hostile reviewers instructed to refute it by finding prior
work. What survives that process is reported here, including the claims that
did not survive.
\end{minipage}

\vfill
{\footnotesize\color{rulegrey}
Companion to \textbf{Mytheus: A Literature Foundation}.\par
Paper numbers cited here refer to the numbered entries in that document.\par}
\vspace{1cm}
\end{titlepage}

\section*{Why this document exists}
\addcontentsline{toc}{section}{Why this document exists}

A final-year project becomes a publishable one at the point where it can state
what it does that the literature does not, and defend that statement against
someone who wants to reject it. Most project reports never attempt this. They
present a system, cite adjacent work approvingly, and leave the reader to
guess at the delta.

This document attempts it directly, and the method matters because it is
adversarial rather than confirmatory. Confirmatory literature review finds
papers that support a design and stops. It reliably produces a claim that
collapses the first time a reviewer asks whether anyone else has done this.
The process behind this document inverted that.

\begin{enumerate}
\item \textbf{Collect the negative space.} Every entry in the companion
document ends with three to five limitations, written to the question ``what
does this paper never test'' rather than ``what did the authors concede''.
That yielded one hundred and thirty two specific failures of coverage.

\item \textbf{Read them through six lenses.} Six independent analyses read all
twenty five limitation sets, each looking for a different kind of absence:
live streaming input, per-recipient output, bounded coordination, evaluation,
grounding and calibration, and systems and deployment. A gap only counts when
several papers fail in the same direction.

\item \textbf{Force the gaps into falsifiable claims.} Candidate contributions
were required to name an experiment, a baseline and a metric. A claim that
cannot be measured was not allowed through.

\item \textbf{Attack every claim.} Each candidate was given to two hostile
reviewers, one searching the academic literature and one searching industry
and adjacent fields such as groupware and information filtering, both
instructed to default to refuted and to cite only prior work they could find
live. Their job was to kill the claim.
\end{enumerate}

The result of step four is the most valuable part of this document, and it is
reported honestly. Where a claim was refuted, it is marked refuted and the
narrower surviving version is presented instead. A contribution that has
already survived a hostile search is worth considerably more than one that has
never been tested, and the narrower claim that survives is usually the one
worth building the paper around.

\vspace{0.4em}
\begin{takeaway}
\textbf{How to use this document.} Read the gap themes to understand the shape
of the opening. Read the contributions section to choose which claim to build
the paper around, paying particular attention to the adversarial review
paragraph in each, since that is what a reviewer will independently reproduce.
Then use the evaluation plan as the experimental programme, because a claim
without its ablation is not yet a contribution.
\end{takeaway}

\newpage
